\section{Additional material for \texorpdfstring{Section~\ref{sec:rel}}{Section 17}: related work}\label{sup:rel}

This section holds the extended discussion of \cref{sec:rel} that the main text summarizes.

\subsection{Formation status}\label{sup:rel-formation-status}

Formation is where the simulations surveyed in \cref{sec:rel-nested} and our results meet. For two groups whose private frames differ by a rotation (\cref{def:frm-two-group}), we prove for the frozen-calibration mean-field map that, above a threshold in policy precision, grounded learning, calibrated precision and attention at finite capacity drive the population from a declared start to the split branch (\cref{thm:frm-formation}). At the exact split, reached with top-$\kappa$ attention ($\gamma=\infty$) and exact in-group translations ($\sigma_\alpha^2=0$), the two groups are exactly the collectives at every level of the window of \cref{thm:frm-split-contrast}(c), a window of admissible levels that is nonempty at every point of the parameter grid we computed (\cref{obs:frm-grid}); at finite policy precision their collective status is a numerical observation at one parameter point (\cref{obs:frm-exploration}). Roles are not supplied as priors, but the split follows a supplied difference of frames, and the result is qualified: the step from the stochastic model to the map is open (\cref{open:frm-closure}); calibration learned online from a symmetric start reaches the split's attention pattern (not necessarily a split selected by contrast) in 0 to 7 of 8 replicates per parameter cell, counts that differ by up to three between reruns of a cell, and is open (\cref{obs:frm-online}); and the collectives arise as a calibration-level effect of misreading joint structure and, at high policy precision, coincide with the partition by codes. From a start law exchangeable given the world path they fail the translation-aware form of the exact-fusion persistence clause (\cref{prop:frm-exact-fusion}), while for full cross-covariance-aware exact fusion this conclusion is heuristic. Persistence with drifting frames (\cref{conj:frm-drifting-frames}) and formation beyond this model (\cref{conj:frm-formation}) are conjectures.

\subsection{Communication models}\label{sup:rel-communication-models}

In the duet of \citet{friston2015duet}, agents with a shared model synchronize, and in its sequel \citep{friston2015hermeneutics} they learn one scalar map from narrative to articulation. Synchrony there comes from shared models rather than from a preference, so the preference to predict and be predicted that we take from the duet needs a formalization of its own; we give two, and at correct codes only being predicted contains coded agreement (\cref{prop:fw-understood-predicted}). Agents that choose partners by expected free energy reach consensus, polarization or splits \citep{albarracin2022epistemic,kastel2023small}. In federated inference \citep{friston2024federated}, which we read in primary text, agents with a common generative model and an assumed common allocentric frame broadcast their posterior predictive beliefs about the shared states, and a receiver reads each broadcast as an observation through a likelihood from beliefs to words. With an identity likelihood this adds the senders' log posteriors to the receiver's own likelihood and prior messages with unit weights (their Eqs.~6 and~8), an update the authors distinguish from Bayesian belief updating. Speaking is a readout of the posterior, and only the direction of gaze is chosen by expected free energy. The likelihood is learned, and in their simulations a common code emerges among naive agents under learning with Bayesian model reduction; language is presented as emerging in a shared econiche, the closest precedent we found for anchoring meaning in a shared world. Reading a partner's belief as an observation is literal reading in our terms (\cref{def:fw-levels}), here with the sender's stated confidence taken at face value, as at level 1. An agent there does not hear itself, which the authors say precludes double counting of its own log posterior; they do not discuss evidence that returns through other agents. \citet{catal2025belief} show that reading beliefs as observations produces echo chambers, which sharing likelihoods removes; this is the double counting of H2 in a discrete model. Apart from \citet{friston2024federated}, we know these works through their abstracts and secondary accounts. The deontic cues of \citet{constant2019regimes} are records in the sense of our record model (\cref{def:fw-record-model}), one mode of a transfer kernel that also covers direct and performative channels (\cref{prop:fw-transfer-kernel}); the word stigmergy is Grass\'e's \citep{grasse1959reconstruction}.

\subsection{Lieb--Robinson bounds and clustering}\label{sec:rel-lieb}
% src: L6-F13, L6-F14, L6 Section 3.1, L4 Sections 2, 3 and 5, L4-P9, L4-P14, B3-T1, B3-T2, B3-R1, B3-P1, B6-F2, B6-F6, B6-F7, decision 16

Hypothesis H4 rests on bounds for the speed of influence and the decay of correlations. \citet{lieb1972finite} proved that finite-range quantum spin dynamics have a finite group velocity, with classical analogs on lattices \citep{marchioro1978velocity,raz2009estimating}. A Lieb--Robinson bound with a spectral gap gives exponential clustering \citep{nachtergaele2006lieb,hastings2006spectral}, \citet{hastings2004locality} treats Markov dynamics on networks, correlations are created only at finite speed \citep{bravyi2006lieb}, and power-law interactions have a hierarchy of linear light cones \citep{tran2020hierarchy}. Our light cone (\cref{thm:sca-lightcone}) is the discrete finite-range case, in which the cone is exact: from a product law, agents more than $2t$ hops apart are independent given the world after $t$ rounds under any local rule, the exact counterpart of the statement of \citet{bravyi2006lieb}. Our clustering theorem (\cref{thm:sca-clustering}) is the classical linear analog of exponential clustering, with the contraction from grounding and sensing gain in the role of the gap. These are known forms applied to a new object, and we claim no novelty for them as mathematics. We change the hypotheses: every length, time and volume statement assumes a stationary world-conditional law with finite second moments and a named contraction (\cref{asm:sca-standing}), because attention selected by expected free energy with sender-specific codes can amplify; without the moment bound only the finite-time light-cone statements survive (\cref{thm:sca-no-moment}), and we call amplifying profiles rumors (\cref{rem:sca-rumors}). For link weights decaying as a power of distance, influence fronts have finite speed exactly when the power is at least $D+1$, the threshold of the free-particle cone of \citet{tran2020hierarchy} (\cref{prop:sca-fronts}).

On $\Z^D$ the covariance decays at least at the Ornstein--Zernike rate \citep{ornstein1914accidental} (\cref{thm:sca-oz}): along an axis the rate $m$ solves $\cosh m=1+\varepsilon/(2a)$, so for a small leak $\xi\approx\sqrt{a/\varepsilon}$ hops, a physical length $\ell\sqrt{a/\varepsilon}$. That the rate is exact pointwise, with the power-law prefactor of the full Ornstein--Zernike form, rests on Green-function asymptotics \citep{ney1966martin} that we know through a restatement \citep{ballu2024asymptotics}, and the constant of the prefactor on a heuristic step, so we state both only as a remark (\cref{rem:sca-oz-form}); the rigorous theory for Ising models above the critical temperature is that of \citet{campanino2003ornstein}. For weights dominated by a reversible chain we prove the form of the Carne--Varopoulos estimate that we use \citep{carne1985transmutation,varopoulos1985long} (\cref{thm:sca-diffusive}). For precisions, inverse-closedness of localized matrices gives some positive decay rate \citep{jaffard1990proprietes}, and our Combes--Thomas bound makes the rate explicit (\cref{thm:col-markov-bound}). In quantum many-body physics the decay length of conditional mutual information diagnoses mixed-state phases \citep{sang2025stability} and is finite for Gibbs states above a threshold temperature \citep{kuwahara2020clustering}; we import it into blanket theory as the Markov length.

The conditional mutual information of interior and exterior paths given blanket paths, the continuous-time counterpart of our window leak (\cref{def:col-blankets}), is proposed as the measure of an inexact boundary by \citet{sakthivadivel2026nonequilibrium}, read in primary text, for continuous-time processes. That paper proves that it vanishes exactly at a path-space boundary, that half of it bounds the mean over boundary paths of the squared total variation distance between the conditional path law and the product of its marginals, and that it bounds the relative entropy between the true and the conditionally separated laws of any measurable observation. The separated law there is separated over the full boundary path, so the last statement does not order the leaks of a fine and a coarse blanket (\cref{ex:met-summary-leak}). In the parts we read it gives no bound on the leak in the coefficients of a model, which \cref{prop:met-leak-bound} supplies for coarse blankets of our class; \cref{thm:col-thick-shell} bounds the instantaneous leak through the decay and conditioning of the covariance.

\subsection{Walk sums, belief propagation and track fusion}\label{sec:rel-fusion}
% src: L6-F15, L6 Section 3.2, L1-D3, L1-P10, L1-X3, L1-X4, L2-P14, B4-P7, B4-P8, B4-O2, B6-F1, decision 6

Hypothesis H2 has precise precedents in estimation. Loopy Gaussian belief propagation converges on walk-summable models with correct means and variances that capture only part of the walks \citep{malioutov2006walk}, correct means hold on any topology when it converges \citep{weiss2001correctness}, and belief propagation is exact on trees \citep{pearl1988probabilistic}. In tracking, double counting of common information is classical: \citet{barshalom1986effect} give the fused covariance of two tracks that share process noise and show that ignoring the dependence overstates the gain from fusion, covariance intersection handles unknown correlations \citep{julier1997nondivergent}, and channel filters fuse exactly on trees \citep{grime1994data}. The ladder from literal reading through the extrinsic message of belief propagation to exact fusion is thus known. We place it inside each receiver's generative model, as a learned model of the sender with a calibration class and an attention policy (\cref{def:fw-levels,def:fw-calibration}), and H2 becomes misreading of joint structure: literal receivers overvalue sources in a stated scope, which is double counting (\cref{prop:mis-overvaluation-scope}), and can undervalue them through common-information cancellation and sender-side completion (\cref{ex:mis-cancellation,ex:mis-sender-completion}).

Consensus filters reproduce the central Kalman filter under exact average consensus at every step, uncorrelated sensor noise and a known number of nodes \citep[Theorem~1]{olfatisaber2005distributed}. Consensus on information contributions reaches the union-model posterior geometrically and recovers that theorem as a limit (\cref{prop:agy-consensus}), while literal consensus, in which each member reads its neighbors' posteriors at level 1 with their posted precisions as law headers, counts a common prior once per member on regular graphs (\cref{prop:agy-literal-consensus}). For exact fusion at a finite number of rounds, which we leave open (\cref{open:agy-exact-fusion-consensus}), \citet{mossel2016efficient} state in their abstract that Bayesian agents with Gaussian signals who exchange estimates reach the full-information belief within $2N\operatorname{diam}(\mathcal G)$ rounds.

\citet{heins2026collective}, which we read in primary text, fits a school-level model to a simulated fish school whose members change speed with local light and turn in response to their neighbors, without estimating the gradient themselves. Patch estimates of the darkness gradient are pooled as natural parameters with the prior counted once (their Eq.~10), the union-model rule of \cref{prop:agy-pooling}(a) when patch evidence is independent. In a static control at 64 agents the ratio of squared error to claimed variance of the pooled posterior is about 16 with overlapping patches and about 10 when each agent is assigned to one patch, and the author attributes part of it to gradients that vary across the school, a misspecification distinct from double counting. The implemented filter adds complete patch posteriors, a product of posteriors in the sense of \cref{prop:agy-pooling}(b). These are simulations of an observer's estimator, with no receiver and no theorem.

\subsection{Strong data processing and broadcasting}\label{sec:rel-sdpi}
% src: L6-F16, L6-P5, L6-C3, L3 Section 1 literature paragraph, L3-P2 to L3-P4, L4-P15, L4-T2, L4-T3, L4-C1, B1-P6, B6-F5

Fidelity is governed by strong data-processing inequalities, which originate with \citet{ahlswede1976spreading}. \citet{polyanskiy2017strong} bound information flow in Bayesian networks by site percolation, following \citet{evans1999signal}, and \citet{polyanskiy2019application} bound the mutual information of distant variables of a Markov random field by bond percolation. For a Gaussian source and an additive Gaussian channel the Kullback--Leibler and $\chi^2$ contraction coefficients coincide \citep{makur2020comparison}, so our coefficient, the top squared canonical correlation (\cref{prop:coh-sdpi}), is known in substance; we prove it directly and derive its consequences in the framework, namely products along relay paths (\cref{prop:coh-relay}), decay of information at twice the mass (\cref{prop:coh-info-decay}) and percolation bounds on the information that a protocol carries about a distant source (\cref{thm:sca-percolation}).

On trees, the binary symmetric channel is reconstructible exactly when its squared correlation exceeds the percolation threshold \citep{bleher1995purity,evans2000broadcasting}, and the Kesten--Stigum condition \citep{kesten1966additional} is sufficient in general but not always sharp \citep{mossel2001reconstruction}. For Gaussian broadcasting on $b$-ary trees with edge correlation $\alpha$ we prove that $b\alpha^2>1$ is necessary and sufficient, with the branching number deciding on general trees away from the critical case by the theorem of \citet{lyons1990random}, read in a review (\cref{thm:sca-trees,thm:sca-gaussian-percolation}). At depth two or more the regime $b\alpha^2>1$ lies, in plain coordinates, outside the diffusive class on which contrast is defined (\cref{prop:sca-tree-diffusive}), and condensation of coherent information on general networks is a conjecture (\cref{conj:sca-condensation}). With power-constrained Gaussian relays, information decays only like $\log\log h/\log h$ over $h$ hops \citep{polyanskiy2016dissipation}; whether relays that re-encode content, such as writing, extend a modality's reach beyond the linear correlation length is open.

\subsection{Group synchronization}\label{sec:rel-sync}
% src: L6-F17, L6-P7, L6-C2, L3-T3, L3-P12, L3-P13, L4-T5, L4-T6, L4-C2, L5-P11, B5-T2, B5-R2, B6-F11 to B6-F13, decisions 9 and 18

Consistency is a synchronization problem. \citet{singer2011angular} recovers angles from noisy offsets spectrally, \citet{bandeira2013cheeger} relate the smallest eigenvalues of the connection Laplacian to the least frustration of an assignment by inequalities of Cheeger type, and \citet{abbe2018group} show that on $\Z^D$ distant group differences can be recovered at small noise for $D\ge3$ but not for some continuous groups at $D=2$. Frustration is the product of signs around a plaquette \citep{toulouse1977theory}, balance of signed and gain graphs is the same condition on cycles \citep{harary1953notion,zaslavsky1989biased}, and with rank-limited maps agreement is a global section of a discourse sheaf \citep{hansen2021opinion}. Our holonomy criterion (\cref{thm:coh-holonomy}), that consensus with a given content exists exactly when the cycle holonomies fix it, imports its form for directed closed walks on strongly connected graphs from the author's gauge-theoretic manuscript \citep{dennis2026magent} and states it on the fundamental cycles of weakly connected graphs, where requiring every content to be fixed gives the balance test of gain graphs. We expect this reduction to fundamental cycles to be classical in the theory of local systems on graphs and claim no novelty for it as mathematics. What the framework adds is the origin of the maps: translations are learned from grounded co-occurrence and true translations are flat (\cref{prop:coh-flat-true}), so holonomy diagnoses errors in learned codes, and frustration acts like grounding (\cref{prop:coh-frustration-grounding}). Whether learned translations inherit the dimension dependence that \citet{abbe2018group} prove for independent noise is open. Cycle amplification under heterogeneous priors is a calibration failure, distinct from holonomy and from double counting (\cref{thm:coh-amplification}), and the absence of cyclic arbitrage says that log exchange rates are a gradient, as in the combinatorial Hodge decomposition used for ranking \citep{jiang2011statistical} (\cref{prop:abs-exchange-consistency}). Noiseless relays of fixed codes in a common metric are alternating projections, which converge to the projection onto the intersection \citep{vonneumann1950functional,halperin1962product} at rates studied by \citet{kayalar1988error} and \citet{badea2012rate}; we prove the rate we use directly (\cref{thm:sca-alternating}), and random hop maps suggest that the rank of shared content falls with scale, generically to one (\cref{conj:sca-rank-down}); the fall need not be monotone when the codes share no direction (\cref{ex:ren-rank-nonmonotone}).

\subsection{Mori--Zwanzig and closure}\label{sec:rel-closure}
% src: L6-F18, L6-P6, L2 Section 4 literature sentences, L2-T4, L2-T5, L7-G15, B3-T3, B3-X1, B3-R3, B4-D5, B4-T3, decision 15; H1-R8, H1-R22 to H1-R24, H1-P3, H1 change list (sec:rel-closure; prop:col-nesting and thm:col-bridge precedents), reading levels raised by H1: barnett2023dynamical to P (Section 2), rosas2024software to P-qfe (main text)

Closure asks when a collective's coarse variables have dynamics of their own. \citet{zwanzig1961memory} and \citet{mori1965transport} derived exact equations for selected variables with memory kernels and fluctuating forces, minimal realization goes back to \citet{ho1966effective}, and aggregation to the nearly decomposable systems of \citet{simon1961aggregation} and the lumpable chains of \citet{kemeny1960finite}. Information theory offers related order parameters: informational closure, which implies sufficiency under a Markovian system--environment structure \citep{krakauer2020information}, causal blankets \citep{rosas2020causal}, which we know through secondary accounts, and dynamical independence \citep{barnett2023dynamical}, whose definition we read in the primary text. Our elimination of a collective's interior, with its memory kernel, its Hankel-rank count of memory variables and its condition for vanishing memory (\cref{thm:agy-mori-zwanzig,thm:agy-closure-identity,rem:agy-minimal-memory}), is linear Mori--Zwanzig in discrete time. What we add is a declared closure property and its link to contrast. A small memory residual, statistical sufficiency and a slow retained spectrum are distinct conditions \citep{dennis2026participatory}; we require the first together with inferential closure (\cref{def:agy-closure}). Contrast controls the memory of the collective average for reversible drift, and a drift-based contrast with row and column dispersions does so for non-reversible drift (\cref{thm:agy-closure-average}); the symmetric dispersion alone does not suffice (\cref{ex:agy-symmetric-dispersion}). The drift-based rates differ from the corresponding Rayleigh quotients of the stationary precision by at most a term set by the stationary probability current (\cref{prop:agy-drift-precision}), and the drift-based inward rate also keeps the couplings to outsiders, so the drift-based contrast can exceed the selecting contrast with no current at all (\cref{sup:agy-memory-and-closure-of-the-average}).

\citet{barnett2023dynamical} call a macroscopic variable dynamically independent of the microscopic system, in the context of an environment, when a conditional transfer entropy from the microscopic past to the macroscopic future vanishes; they compute it for linear state-space systems and identify the condition with the informational closure of \citet{bertschinger2006information}, which we know through their restatement. \citet{rosas2024software}, whose main text we read through targeted extraction, define informational, causal and computational closure for stationary processes in discrete time, prove the first two equivalent for macroscopic variables that are functions of the microscopic process, and characterize closed levels by strong lumpability of causal states, which are maximally predictive and minimal \citep{shalizi2001computational} and, according to an abstract, are macrostates \citep{shalizi2025macrostate}; the framework has no environment, no blanket and no bound for approximate closure. For linear-Gaussian collectives with full-rank noise two of these notions, dynamical independence and informational closure in its one-step form, meet ours (\cref{sec:met,prop:met-closure-notions}): in the stationary case without input the dynamical dependence of a coarse variable is at most the closure defect of \cref{def:agy-closure}, and in that case the exact forms of dynamical independence, one-step informational closure and the closure of \cref{thm:agy-closure-identity}(b) coincide.

Near decomposability, known to us from an abstract \citep{simon1961aggregation}, a secondary account \citep{simon1962architecture} and a record \citep{courtois1977decomposability}, is the picture behind the nesting of spectral contrast at weak grounding, through a hierarchy of timescales (\cref{prop:col-nesting}; \cref{thm:col-bridge}(iv)). Fast variables are eliminated with corrections in the ratio of timescales \citep{vankampen1985elimination,geigenmuller1983systematic,haken1982slaving} and rigorously in the limit of infinite separation \citep{pavliotis2008multiscale,berglund2006noise,givon2004extracting}, works we know from records and abstracts. Whether any of them bounds the memory of coarse variables of our kind at a finite ratio we have not checked; \cref{thm:agy-closure-average} and \cref{sec:met,sec:ren} give such bounds through contrast.

\subsection{Potential games and logit dynamics}\label{sec:rel-games}
% src: L2-P15, L2-T6, L2-P16, L2 Section 6 literature sentences, B4-L1, B4-T1, B4-T2, B4-X1, B5-P10 (scope as corrected in Section 6), B6-F9, B6-F10, L6-O7, L7-M32, B2-O1

Agency asks when members' separate choices act as the choices of one agent. A game is an exact potential game if and only if every four-cycle of profiles has zero total payoff change \citep[Theorem~2.8]{monderer1996potential}, a statement we read in restatements \citep{arefizadeh2024characterizations}. Asynchronous logit revision then has a Gibbs stationary law \citep{blume1993statistical}, which we know secondarily, its convergence depends on asynchrony \citep{alosferrer2010logit}, and near-potential games have logit laws near a Gibbs law \citep{candogan2013dynamics}. \citet{wolpert2001optimal} design members' payoffs so that individual optimization serves a world utility. Our agency theorem (\cref{thm:agy-agency}) finds such a potential instead of designing it: under its protocol, in which members share belief and data, announce their policies, share one precision and deliberate asynchronously, each member's expected free energy has the group expected free energy as an exact potential, while synchronous deliberation of two members with redundant slots does not reproduce the group posterior (\cref{prop:agy-synchronous}). With heterogeneous beliefs the deliberation law stays within explicit distances of a Gibbs law of a potential (\cref{thm:agy-heterogeneous}), with a bound on the ratio of transition probabilities sharper than that of \citet{candogan2013dynamics}, and joint acts keep a potential when members agree on how their acts combine and lose it when interventions interact and the members' displacements, predicted outcome minus target, differ along a joint effect (\cref{thm:agy-joint-acts,ex:agy-no-potential}). The code game with a common value is an exact potential game on rings and on complete graphs, and on open chains of at least three agents only in its pairwise form (\cref{prop:coh-code-potential}). The stochastic-approximation step behind formation, of the kind treated by \citet{benaim1999dynamics}, is not carried out (\cref{sec:frm}).

\subsection{Opinion dynamics and social learning}\label{sec:rel-opinion}
% src: L6-F19, L6-P8 (Friedkin-Johnsen reading; partial-fraction form of thm:sca-oz(a)), L6-F15 (social learning), L3 Section 2 literature paragraph, L3-P8, L3-T2, L1-P15, L1-C1, B4-P8

Repeated weighted averaging \citep{degroot1974reaching}, known to us secondarily, and the influence process of \citet{friedkin1990social}, whose equilibrium is a resolvent of the influence matrix with DeGroot's model as a limit, are the basic linear models of opinion dynamics; bounded confidence fragments opinions \citep{hegselmann2002opinion,deffuant2000mixing}, a rule that senders who want to be predicted follow on the sender's side (\cref{prop:fw-bounded-confidence,conj:fw-fragmentation}); and in noisy consensus the variance of deviations from the average, which \citet{bamieh2012coherence} call coherence, scales badly with network size in one and two dimensions. Our ring, with $W_{ii}=c$, $W_{ij}=a$ between neighbors and local noise $qI$, carries a Friedkin--Johnsen influence structure with grounding and forgetting as stubbornness. Its world-conditional covariance $q(I-W^2)^{-1}=\tfrac q2[(I-W)^{-1}+(I+W)^{-1}]$ decays at the rate $m$ of the total-influence matrix $(I-W)^{-1}$, $\cosh m=1+\varepsilon/(2a)$, because the second term decays at least as fast, at $\cosh m'=(1+c)/(2a)$. So the physical correlation length $\ell\xi\approx\ell\sqrt{a/\varepsilon}$ is the Friedkin--Johnsen influence length (\cref{thm:sca-oz}), and its square grows like the diffusion constant $a\ell^2$ times the correlation time (\cref{prop:sca-diffusive-scaling}). Approximately exponential decay of influence has been reported for forced linear diffusion models \citep{murraywatters2026decay}, Friedkin--Johnsen dynamics have been written with a Green operator \citep{boudourides2026boundary}, and $\ell\sqrt{a/\varepsilon}$ is the decay length of the standard massive lattice Green's function, so we do not present it as new. The result of \citet{bamieh2012coherence} is the control-theoretic counterpart of our massless limit (\cref{thm:coh-grounding-tension}).

Social learning names the overvaluation case of H2 several times: persuasion bias \citep{demarzo2003persuasion}, naive herding \citep{eyster2010naive} and correlation neglect \citep{enke2019correlation}, with conditions for the wisdom of DeGroot learners \citep{golub2010naive}. \citet{madsen2018large} report echo chambers among Bayesian agents who treat received beliefs as fresh evidence and drop partners outside a confidence window; in our terms these are literal receivers with withdrawal by disagreement, consistent with our results that weak and self-sealed echo chambers require misspecification (\cref{thm:mis-echo-chambers,prop:mis-exact-no-echo}). \citet{nguyen2020echo} separates epistemic bubbles, which omit outside voices, from echo chambers, which discredit them; in our terms these are closure of attention (\cref{def:mis-echo}) and miscalibrated precision.

\subsection{Disorder and convention domains}\label{sec:rel-disorder}
% src: L6-F20, L6-C1, L3-T5, L3-C2, L5-L2, L5-P10, B5-P9 to B5-P11, B5-T1, B5-C1 to B5-C3, B5-C2-lit, C2 S5, repair coh(i); ring scope of the code game per the Section 6 correction of B5-P10

\citet{imry1975random} argue that a weak random field destroys the order of systems with continuous symmetry below four dimensions, by the domain argument that bears their name; \citet{aizenman1989rounding,aizenman1990rounding} prove that quenched randomness suppresses continuous symmetry breaking up to four dimensions and leaves the two-dimensional random-field Ising model with a unique Gibbs state, whereas the three-dimensional model has a phase transition \citep{bricmont1988phase}. For the two-dimensional model, \citet{ding2023correlation} prove that the box size at which the influences of the field and of the boundary condition on the magnetization are comparable grows, at zero temperature, exponentially in the power $-4/3$ of the field strength relative to the coupling. Dialect variation grows sublinearly with distance in the studies since \citet{seguy1971relation} that \citet{nerbonne2010measuring} reports, and \citet{burridge2017spatial} models dialect boundaries as interfaces with surface tension produced by a conformity nonlinearity.

We derive that conformity: under calibrated reception and uniform habits, expected free energy gives a logistic adoption rule, conformist above a threshold in policy precision for an affine value without pull, with an innovation floor that acts as a leak; at equal frequencies it has the value and slope of the power conformity rule of dialect models at matched exponent, which has no such floor (\cref{prop:coh-adoption,prop:abs-conformity}). Because the code game with a common value is an exact potential game on rings (\cref{prop:coh-code-potential}), the one-dimensional domain law, with the Imry--Ma scaling, is a theorem there (\cref{thm:coh-domains-1d}); on open chains the same bound holds for the pairwise game. In higher dimension the code game is an exact potential game only for special value functions, and coarsening, domain sizes under random pulls and the bundling of isoglosses remain conjectures (\cref{conj:coh-coarsening,conj:coh-imry-ma,conj:coh-bundling}); the two-dimensional conjecture may carry the exponent $4/3$ if domain size is identified with the length of \citet{ding2023correlation}, a heuristic of ours and not their theorem. Sparse, unclustered walls make S\'eguy curves rise linearly with distance, against the square-root rise of a leaky field with exponential correlation (\cref{prop:abs-seguy}).

\subsection{Abstract collectives in the social sciences}\label{sec:rel-social}
% src: L5-L1 to L5-L15, L5 Sections 2 to 6, L6-F21, B5-D3, B5-L2, B5-P12, B1-T3(iii), decision 12

\citet{deutsch1953nationalism} defined a people by the complementarity of its members' communication, tested by how fast and accurately messages get through; we read the defining chapter in the 1966 edition. \citet{savage1960statistical} compare transaction flows with a null model built from marginal totals, and modularity compares within-group links with a null model \citep{newman2006modularity}. These are the closest precedents for contrast and for declared nulls (\cref{def:col-contrast,def:agy-nulls}); the transaction-flow index measures contrast in message rates rather than in coherence, and our barrier and bypass results connect the two (\cref{prop:abs-barrier,prop:abs-bypass}). The ultrametric contrast of societies uses world-conditional correlations, since under a common mode unconditional ones cap it (\cref{def:col-ultrametric,prop:abs-ultrametric-cap}).

\citet{hayek1945use} held that dispersed knowledge is communicated by prices, and rational-expectations aggregation \citep{grossman1976efficiency,hellwig1980aggregation,grossman1980impossibility} is cited for attribution only. A Walrasian price formed from private signals counts the common prior once per trader, while a revealing price is sufficient (\cref{prop:abs-walras}), and public prices suffice exactly when the decision functionals lie in their row space (\cref{prop:abs-hayek}). Innis's time- and space-biased media \citep{innis1950empire}, known to us secondarily, map onto the decay and reading range of records (\cref{prop:sca-records,prop:sca-records-length}); the mapping is ours.

Borders reduce trade between industrialized countries by 20 to 50 percent \citep{anderson2003gravity}, and the distance effect on trade has persisted \citep{disdier2008puzzling}. Single-scale diffusive coherence in two dimensions predicts that distance elasticities fall as communication lengthens the correlation length, contrary to that persistence, while nested coupling gives scale-free elasticities (\cref{prop:abs-hierarchy,prop:abs-elasticity}). \citet{alesina1997number} trade the gains of size against the costs of heterogeneity; our criterion compares a polity's fixed cost with the heterogeneity weight times the correlation volume of preferences, under two declared policy rules (\cref{prop:abs-polity}). Marchetti's daily travel budget of about one hour \citep{marchetti1994anthropological}, constant only in aggregate \citep{mokhtarian2004ttb}, bounds the diameter of one-round sets (\cref{prop:abs-one-round}), and the territorial-unity condition behind urban scaling \citep{bettencourt2013origins}, known to us through \citet{ribeiro2023mathematical}, is a one-round reach condition with a money budget. We read the scale-free correlations of starling flocks \citep{cavagna2010scale} as the massless regime, where the size bound of \cref{cor:col-lstar} is vacuous. These mappings are interpretations, and the framework does not cover ungrounded conventions, strategic deception, power or non-stationary history (\cref{sec:abs}).

\subsection{Hierarchical active inference}\label{sec:rel-hierarchy}
% src: H1-R10 to H1-R13, H1-P2, H1-R33, H1 Section 9 (H5b; claims 5 and 6), H1 change list (prop:agy-consensus, prop:agy-literal-consensus, prop:mis-believed-values), HC Section 4 (Friston 2008 verified at the level of the phrases, equation numbers not rechecked; Friston et al. 2017 verified), HC V1(ii) (reconciled gloss), HC 2.8 (reading levels), decisions 6, 10 and 20

In hierarchical dynamic models \citep{friston2008hierarchical} the causes of each level are generated by the level above, whose prediction is an empirical prior on the level below, and inversion is a gradient flow on free energy in which precision-weighted prediction errors ascend and predictions descend, the predictive coding of \citet{rao1999predictive}, cited for attribution. \citet{kiebel2008hierarchy}, read through targeted extraction, impose a separation of timescales in which a slow level is a control parameter of a fast one, and state that it is not a property of coupled neuronal systems as such. In the deep temporal models of \citet{friston2017graphical} a higher state generates the initial state of a lower sequence, the ascending message is a Bayesian model average and the descending one an empirical prior, and for continuous states belief propagation and variational message passing reduce to predictive coding; \citet{friston2017deep} and \citet{pezzulo2018hierarchical}, known to us from records, extend the architecture to planning and control.

As we read them, these hierarchies do not double count. Their updates are steps on one free energy of the joint model, in which each factor appears once, and on linear-Gaussian trees predictive coding and mean-field message passing reach the exact posterior means (\cref{prop:hie-predictive-coding}), while belief propagation excludes the recipient's own message by the cavity rule \citep{pearl1988probabilistic}. In the two-pass form of exact inference on trees the upward message is a likelihood, equivalently a prediction error relative to the cavity prior; in predictive coding errors are computed against current estimates, with top-down predictions weighted by the generative precision, and that scheme is also exact. Double counting arises only when a posterior that already contains the receiver's data is used as an independent prior, which is level-1 reading of a message whose provenance overlaps the receiver's data, the overvaluation case of H2 (\cref{def:fw-levels,prop:mis-overvaluation-scope}).

In these hierarchies the level above is another part of one model. We found no treatment in which the top-down prior is built from the receivers' own messages, as when members read a summary of the collective that they wrote, and none that analyzes such a circular hierarchy with interpretation levels, calibration classes, records or a null family that includes exact fusion. \Cref{sec:hie} does so in linear-Gaussian models: a member's two-level model of its collective as a latent cause is correct exactly when the collective coordinate is closed, and literal top-down reading, done once, overcounts by an overconfidence factor below two in a scalar tree, and by factors growing linearly or geometrically in the number of rounds when a persistent record is rebuilt from literal posteriors. That prediction errors between agents cannot detect shared error (\cref{lem:mis-social-blind}) has no counterpart in the hierarchical literature we read. In the author's gauge-theoretic program, as secondary descriptions summarize it, the parent's belief is written into each child's prior with a blend factor \citep{dennis2026formation}, a circular case whose exact top-down message is the cavity message.

\subsection{Cultural priors and top-down social signals}\label{sec:rel-culture}
% src: H1-R20, H1-R21, H1 Section 9 (H5b; claims 1 and 8), H1 change list (prop:coh-adoption, thm:coh-grounding-tension), HC V3, V4, V13, decisions 12, 13 and 22

Cultural priors reach individuals through inference of other agents' expectations in shared practices \citep{veissiere2020thinking}, known to us from its abstract, through conventional affordances that rest on shared expectations rather than on correlations in the environment \citep{ramstead2016cultural}, and through deontic cues that consolidate as repeated actions raise the counts of the environment's likelihood parameters \citep{constant2019regimes}, with niche and organism treated symmetrically by \citet{constant2018variational}. In our terms these priors travel through records (\cref{def:fw-record-model}) and behavior, and none of these works asks whether members' own contributions to the cues are counted twice when the cues are read back. Synergetics, cited here for the slaving principle \citep{haken1983synergetics}, has order parameters that the parts produce and that enslave them, a circularity carried to closed equations for opinion shares by \citet{weidlich1991physics}, and \citet{galam1986majority} renormalizes majority opinion up a hierarchy of committees; we know the last two secondarily. According to its abstract, \citet{morris2002social} find that more precise public information raises welfare when agents have no socially valuable private information and has ambiguous effects otherwise, which bears on the precision balance of \cref{thm:coh-grounding-tension}.

None of these is a linear-Gaussian theory with calibration tested against a null. In \cref{sec:hie} a member-written record that does not decay lengthens the memory of the collective mode and extends identity under turnover, at the price, for a grounded coordinate, of a shared-error floor that is nondecreasing in the weight members give the record; the precision-based contrast computed with the record as a node, at the record noise that makes the joint dynamics reversible, gains a factor independent of the write weight when the record is read strongly enough, and it is undefined for a noiseless global record as soon as two members are uncoupled. Linear conformity to a self-constituted prior selects no partition that the declared couplings and readership do not already contain, and formation by conformity is a conjecture.

\subsection{Renormalization in active inference}\label{sec:rel-renormalization}
% src: H1-R1 to H1-R6, H1-R9, H1-P1, H1-O1, H1-O2, HC Section 4 (the grouping operator forms the partition and eliminates internal states by conditioning on blankets, the reduction operator keeps the slow eigenmodes of each blanket's Jacobian; 32 singular vectors in the main text and at most 16 in the caption of Figure 4 of Friston et al. 2025; Parcels and particles not available to the check, kept at its reading level), H4-R2, H4-P9, H4-T8, HC V10, decisions 16, 19 and 21

\citet{friston2019free} builds scales with two operators, in the parts we read. A grouping operator forms a particular partition from the Jacobian, a partition the text calls one of many, and eliminates internal states by conditioning on the blankets, adding the remainder to the fluctuations; a reduction operator keeps the slow eigenmodes of each blanket's Jacobian and adds the fast ones to the fluctuations. The composite is called a renormalization in the working sense of \citet{lin2017deep}, known to us from its record. In our reading it assumes a sparse Jacobian read as a blanket, a heuristic partition with free thresholds, elimination given the particle's own blanket at one time, white intrinsic fluctuations independent across blankets, a spectral gap in each blanket's Jacobian with the fast modes of other blankets treated as noise, Gaussian higher-level noise by a central-limit argument, and slow and fast modes of different particles that decouple, with noise correlations neglected. Accuracy rests on the separation of timescales, with the slaving principle \citep{haken1983synergetics} and center manifold theory cited, and we found no error bound in the pages we read.

\citet{friston2021parcels}, read through targeted extraction, apply the operators to one subject's imaging data and report time constants growing by a factor of about 2.37 and spatial support by about 2.13 per scale, figures we have not verified, for a partition that they state is not unique. \citet{friston2023path} call their fixed flow operator an adiabatic approximation and list particles of particles as an open question, without a renormalization theorem. \citet{friston2023variational} give a population at a slow scale its own partition into external, sensory, active and internal kinds, the closest precedent we found, in conceptual form, for coarse external states defined as slow variables of other collectives. The renormalizing generative models of \citet{friston2025pixels} build conditional independence into a discrete hierarchy with one parent per state and keep the leading singular variates of image blocks; they decline to benchmark the examples, although they report MNIST classification accuracy (95.1\% on the full test set, 99.8\% on the half with higher marginal likelihood), and, in the parts we read, give no bound on what the reduction loses.

What these constructions lack is an account of what the reduction drops, which \cref{sec:ren} supplies for linear-Gaussian models. The elimination of internal states is the Mori projection of \cref{thm:agy-mori-zwanzig}(b) with its memory kernel, initial slip and colored noise discarded; it is exact for every initial law only when the retained states do not read the internal ones, so never for a particle whose active states read its internal states, and its error is, to leading order, a computed coefficient times the inverse ratio of timescales when the noise factors are invertible, the internal drift and the quasi-static drift are stable and the internal noise intensity is scaled with that ratio (\cref{prop:ren-adiabatic-rate}), as two computed instances, not tested against these hypotheses, show once the ratio is ten or more (\cref{obs:ren-adiabatic-rate}). For the linear reversible analogue of the adiabatic reduction the static error and the memory are bounded by spectral quantities of the partition, and for block averages through the dispersion of the cut over the internal stiffness, which high contrast alone does not make small (\cref{thm:ren-memory}(d), \cref{sec:ren-closure}). That no error bound accompanies the active-inference renormalization is a search result covering the papers above and the replies on blankets, not the whole literature; it does not cover two later works that we have not read, the treatment of renormalizing generative models by \citet{zaghw2026renormalising} and the engineering perspective on active inference of \citet{devries2026active}.

The elimination conditions on a particle's blanket at one time. \citet{virgo2022embracing}, a commentary without a theorem that we read in the authors' preprint, argue by d-separation in a time-unrolled discrete-time sensor-motor loop that internal and external states are independent given the sensorimotor history together with the initial states, or given the infinite history for an ergodic system, and in general not given the current sensory and active states. They treat one loop, not collectives or renormalization; applying the argument to the elimination is our reading. For sets of agents with linear-Gaussian dynamics and per-channel noise, from a start that the blanket screens, history suffices when no sensory channel reads both sides in one round, with a converse outside a set of coefficients of positive codimension (\cref{thm:col-path-blanket}); the instantaneous leak is bounded by the memory of either side, and the boundary's filtered estimate of the exterior restores the instantaneous blanket (\cref{prop:col-memory}).

\subsection{Block spins and network renormalization}\label{sec:rel-network}
% src: H1-R28 to H1-R31, H1 Section 9 (H5d), H1 change list (lem:col-kron, prop:abs-hierarchy), H4 Section 5 (H4-P16, H4-P17, H4-X4), H4 section-ready summary, HC V8, V9, decisions 16 and 22

\citet{kadanoff1966scaling} takes the magnetizations of cells much smaller than the coherence length as collective variables, the block counterpart of our hierarchies ending where a coarse cell reaches the correlation length (\cref{cor:col-lstar,sec:ren}), a connection of ours; the renormalization group \citep{wilson1971renormalization,wilson1974renormalization,cardy1996scaling}, known to us from records and secondary accounts, is cited for the standard theory, and the exponent of the leak under long-range couplings (\cref{prop:ren-exponents}(b)) is the thermal eigenvalue of the long-range Gaussian fixed point of \citet{fisher1972critical}, cited from its record and abstract. \citet{jonalasinio2001renormalization} presents it as a limit theorem with a Gaussian fixed point and shows that the Gaussian block recursion of Dyson's hierarchical model \citep{dyson1969existence} is exact; our nested medium (\cref{prop:abs-hierarchy}) is a massive model of this type and renormalizes exactly (\cref{sec:ren}). Among information-theoretic formulations \citep{mehta2014exact,beny2015information,kochjanusz2018mutual}, \citet{gordon2021relevance} state in their abstract that bottleneck compression keeps the operators of lowest scaling dimension, which ties the selection of sensed variables (\cref{sec:rel-sensing}) to renormalization, without saying whether Gaussian theories on graphs are covered.

For networks without coordinates, Laplacian renormalization defines supernodes through diffusion \citep{villegas2023laplacian}, beside geometric and multiscale renormalization \citep{garciaperez2018multiscale,garuccio2023multiscale,gabrielli2025network}, and, according to the abstract of \citet{dorfler2013kron}, Kron reduction maps a Laplacian with self-loops to another such Laplacian, which is the decimation step of \cref{lem:col-kron}. A diffusive world-conditional precision with positive leaks, which \cref{thm:col-diffusive-precision} provides under stated conditions, is a massive Gaussian free field on the communication graph, where these operations are exact. Under \cref{asm:ren-standing}, \cref{sec:ren} gives exact flows of the axis rate and the second-moment length, shows that block averages in general leave the diffusive class, so that contrast between meta-agents uses a coarse medium (\cref{def:col-coarse-medium}), and finds that Laplacian renormalization does not see a uniform precision leak, agrees with the criterion $r\ge r_0$ for a population formed by two blocks of equal size with an equitable cut and equal internal gap, under the gap condition of \cref{prop:ren-lrg-agreement}, only on the interval of leaks of that proposition, which contains the window of \cref{thm:col-bridge}(iv), and in an example with interleaved scales returns the selected regions at no single diffusion time (\cref{sec:ren-network}).

\subsection{Predictive information, information bottlenecks and rational inattention}\label{sec:rel-sensing}
% src: H1-R26, H1-R27, H1-R17 (the selection principle stated in words), H1 Section 9 (H5c), H1 change list (conj:sca-rank-down), H5-T1, H5-R1 and the H5 section-ready summary as specified for sec:env, HC V11, V12, HC 2.6 (the cell example uses the buoyant mass in sec:env; H1-E1 is not transcribed)

Predictive information, the mutual information between past and future, measures the complexity of a process \citep{bialek2001predictability}, and the information bottleneck keeps what a compressed signal says about a relevance variable \citep{tishby2000information}. For jointly Gaussian variables the optimal bottleneck projects onto canonical-correlation directions that enter at critical trade-offs \citep[Theorem~3.1]{chechik2005information}; the past-future bottleneck of \citet{creutzig2009past}, read in its abstract, adds state dimensions at structural transitions, and \citet{still2014information} treats predictive inference by the same method; and the optimal predictive representation of a driven oscillator rotates as the prediction lag grows \citep{sachdeva2021optimal}, read through targeted extraction. Rational inattention limits the capacity of the channel from observations to actions \citep{sims2003implications}, and price-setting firms attend to firm-specific conditions when these are volatile or consequential \citep{mackowiak2009optimal}, a result we know from its abstract, as we know the review of \citet{mackowiak2023rational}. \citet[Section~7.2]{fields2022competency} state in words that organisms model only the principal components of their experience.

The selective-sensing results of \cref{sec:env} are thus anticipated in substance: for linear-Gaussian exteriors, information-optimal sensors read leading canonical variates between present and future with allocations of the water-filling form, and error-optimal ones read leading eigenvectors of a decision matrix with reverse water-filling. For the joint treatment of a cap on the number of channels and a budget we claim no priority: the budget form is the information curve of \citet{chechik2005information}, whose Section~8 also names a dimension fixed in advance, and \cref{sec:env} records that \citet{scharf1998wiener} treat rank reduction and a rate budget together in canonical coordinates. What may be new is at most the dependence on horizon and control, for which the selection of slow features at one step has a precedent in \citet{creutzig2008predictive}, read in its abstract, together with the remark that the one-step attention expected free energy over world channels is horizon-free (\cref{prop:fw-efe-communication,rem:env-efe-horizon-free}); a selection principle of coupling, predictability and relevance; and conditions under which aggregates of many agents form a closed environment. Rank cascades concern compression, not products of relay maps, and neither prove nor refute \cref{conj:sca-rank-down}. The leading concentration constant of \cref{sec:env} is the known result for a sphere that counts the molecules inside its volume without hindering their diffusion, which \citet{mora2010limits} give as their Eq.~10 and attribute to \citet{berg1977physics}, known to us from its abstract; the appendix rederives it, and only the finite-time correction and the bound on $\mathfrak q_T$ are new.

\subsection{Organizations as information processors}\label{sec:rel-organizations}
% src: H1-R32, H1-R27 (Daft and Weick), H1 change list (prop:abs-walras, prop:abs-hayek), H3-E3 and H5-E2 as scoped by HC V12, M_spec_H5 Section 3 (V12: worked examples, readings labeled as interpretations)

Economics treats the firm as an information processor: team theory with its linear-quadratic-Gaussian case \citep{radner1962team,marschak1972economic}, decentralized processing under delay \citep{radner1993organization}, hierarchies as allocations of knowledge \citep{garicano2000hierarchies}, organizations as costly channels with a code \citep{arrow1974limits,cremer2007language} and design as information processing \citep{galbraith1974organization}; organizations also scan their environment as interpretation systems \citep{daft1984toward}. We know these works from their records, the formal code through a secondary account and the interpretation systems from an abstract, and cite them for attribution, as precedents for a firm whose environment is a market and for low-rank market signals (\cref{prop:abs-walras,prop:abs-hayek}). This supplement has worked examples in linear-Gaussian models (\cref{ex:met-two-firms,ex:env-firm}), two firms trading on one price and a firm that reads a futures price at its production lag, whose readings for corporations are interpretations, and whether the linear-quadratic-Gaussian teams of \citet{radner1962team} contain such statements we have not checked; the sufficiency of the futures price at the firm's own lag is a known result recovered, the separation theorem of production under price uncertainty in its risk-neutral form \citep{feder1980futures}, which we know from its record and abstract.
